\begin{abstract}
Gradient descent updates the weights of a factored network, but what the network learns is governed by how the \emph{operator} it computes moves: the input-to-output Jacobian $P(x)$, defined by $f(x) = P(x)\,x$ for a bias-free MLP. For deep linear networks this induced motion is known exactly --- under balanced initialisation the gain on singular direction $k$ is $L\,s_k^{2-2/L}$, and the low-rank bias, sequential mode learning and saddle-to-saddle dynamics all follow from that single exponent. Balancedness does not hold at the initialisations used in practice, however, and no analogue of the law is known once the gates depend on the input; the results that do address nonlinear networks concern endpoints, infinite-depth limits or ranks rather than finite-time rates.

We introduce a diagnostic that measures this induced motion in any factored network --- deep linear and MLP alike --- and that reproduces the balanced closed form to machine precision. The measurement admits an exact decomposition: the gain exponent equals $2$ plus the rate at which a singular-vector misalignment between the operator and its layer-wise contexts grows as the singular value shrinks. Under balanced initialisation that rate is exactly $-2/L$, which is precisely what produces $2-2/L$; away from balance it is steeper, and the shortfall in the exponent is exactly the excess. We further show that realising a prescribed operator step in
the weights is an exactly solvable problem so that training may be driven either by gradient descent's own step or by the best reachable one.

This apparatus is built to settle three questions. Does the gain law survive input-dependent gates, and with what exponent? Is the exponent a property of the architecture or of the update rule, a question that cannot be settled from gradient-descent trajectories alone and that the solvable step is designed to answer? And when does a difference in the training dynamics become a difference in what is finally learned?
\end{abstract}
